% !Mode:: "TeX:UTF-8"
\chapter{相关理论与基础技术}
\label{chap:theory} 

\section{点云数据基础}

\subsection{点云数据定义与特性}
点云是三维空间中离散采样点的集合，是三维物体与场景最常用的数字化表示形式之一，其核心作用是通过离散的采样点，精准还原现实世界中物体的三维几何形状与表面属性。从数学定义上来看，一个完整的点云可以表示为 $\mathcal{P} = \{p_i \in \mathbb{R}^3 | i=1,2,...,N\}$，其中 $N$ 为点云中点的数量，$p_i$ 代表点云中的第 $i$ 个点，其三维坐标 $(x_i,y_i,z_i)$ 是点云最核心的基础信息。除了三维坐标之外，根据采集设备的不同，点云中的每个点还可以包含额外的属性信息，包括颜色信息、法向量、反射强度、纹理坐标等，这些属性信息能够为后续的点云处理任务提供更多的有效信息，提升特征提取的精度。

与二维图像规则的网格结构不同，点云数据具有其独特的特性，这些特性也决定了点云处理任务的难度与特殊性，是设计点云处理算法时必须考虑的核心因素。

首先是点云的无序性，点云中的点是三维空间中的离散采样点，其排列顺序不具备任何语义信息，任意改变点云中各个点的排列顺序，都不会改变点云所代表的三维几何形状。这一特性使得传统基于规则网格结构的CNN架构无法直接应用于点云处理，因为CNN的卷积操作对输入数据的排列顺序具有强依赖性，无序的点云输入会导致CNN提取的特征出现严重的歧义，无法有效表征点云的几何结构。因此，点云处理算法必须具备排列不变性，无论点云中的点如何排列，都能提取到一致的特征表示。

其次是点云的稀疏性与非均匀性，在实际采集过程中，受设备性能、采集距离、物体表面反射特性等因素的影响，点云中的点在三维空间中的分布往往是不均匀的。对于距离采集设备较近的区域，点云密度较高，而距离较远的区域，点云则较为稀疏；同时物体的平坦表面区域点云分布均匀，而边缘、角落等细节区域，往往会出现采样不足的问题。这种稀疏性与非均匀性，给点云的局部特征提取带来了极大的难度，算法需要能够适配不同密度的点云分布，在稀疏区域也能准确提取有效的几何特征。

第三是点云的不规则性，现实世界中的物体几何结构千差万别，对应的点云在三维空间中的分布也没有固定的规则，不存在二维图像中固定的邻域关系。对于二维图像而言，每个像素的邻域是固定的上下左右相邻像素，而点云中每个点的邻域，需要根据点的空间位置动态构建，不同区域的邻域点分布差异极大。这种不规则性，要求点云处理算法能够动态构建点的邻域关系，有效提取局部几何特征，同时建模局部区域之间的全局关联。

这些特性共同决定了点云处理任务需要设计专门的网络架构与处理方法，以适配点云的独特属性，这也是点云补全任务的核心难点所在\cite{xiang2021snowflakenet}。

\subsection{点云预处理方法}
在实际的点云处理任务中，原始采集得到的点云数据往往包含大量的噪声、离群点、冗余信息，同时存在尺度不一、坐标偏移等问题，无法直接输入到网络中进行处理。需要通过一系列的预处理操作，提升点云数据的质量，消除无关干扰信息，使其符合网络的输入要求。常用的点云预处理操作主要包括去噪、下采样、归一化、法向量估计等，这些操作是点云处理全流程的基础，直接影响后续任务的精度与效果。

点云去噪的核心目标是去除原始点云中的噪声点与离群点，这些异常点主要来自于采集设备的测量误差、环境光干扰、物体表面的高反光特性等，会严重破坏点云的几何结构，影响后续的特征提取与处理。常用的去噪方法包括统计滤波与半径滤波，统计滤波的核心原理是对每个点，计算其到指定邻域内所有点的平均距离，假设所有点的平均距离符合高斯分布，通过计算分布的均值与标准差，将平均距离超出设定阈值范围的点判定为离群点并去除。该方法能够有效去除分布较为分散的离群点，是实际应用中最常用的去噪方法之一。半径滤波的核心原理是对每个点，统计其在指定半径范围内的邻域点数量，若邻域点数量少于设定的阈值，则将该点判定为离群点并去除。该方法能够有效去除局部的孤立噪声点，计算效率更高，但对滤波半径的设置较为敏感。

点云下采样的核心目标是在尽可能保留点云核心几何结构的前提下，减少点云中点的数量，降低后续处理的计算量，同时让点云的分布更加均匀。原始采集的点云往往存在密度不均的问题，部分区域点云密度过高，存在大量的冗余信息，会增加不必要的计算量，同时也会影响网络的特征提取效果。常用的下采样方法包括体素下采样与最远点采样，体素下采样的核心原理是将点云所在的三维空间划分为固定尺寸的三维体素网格，对于每个体素网格内的所有点，计算其重心点，用该重心点代表体素网格内的所有点，实现点云的下采样。该方法计算效率高，能够快速得到分布均匀的下采样点云，缺点是会改变原始点云的坐标，对于细节区域的几何结构会有一定的损失。最远点采样的核心原理是通过迭代的方式，从点云中选择距离已选点集最远的点，逐步生成下采样点集。该方法能够在指定采样点数的前提下，尽可能均匀地覆盖整个点云，最大程度保留点云的几何结构，不会改变原始点的坐标，是点云深度学习任务中最常用的下采样方法，但其计算复杂度高于体素下采样。

点云归一化的核心目标是消除不同点云之间的坐标偏移与尺度差异，让所有输入点云都处于统一的坐标空间与尺度范围内，提升网络的训练稳定性与泛化能力。常用的归一化操作包括中心化与尺度归一化，中心化的核心原理是计算点云所有点的坐标均值，将每个点的坐标都减去该均值，让点云的重心移动到坐标原点，消除坐标偏移带来的影响。尺度归一化的核心原理是计算中心化后的点云所有点到原点的最远距离，将每个点的坐标都除以该最远距离，让点云被归一化到单位球体内，消除不同物体之间的尺度差异，确保所有输入点云都处于相同的尺度范围内。

法向量估计是点云预处理中的重要操作，点云的法向量信息能够反映物体表面的局部几何朝向，是后续三维表面重建任务的核心基础信息，同时也能为点云特征提取提供额外的几何信息。常用的法向量估计方法是基于主成分分析（Principal Component Analysis, PCA）的方法，其核心原理是对于每个点，选取其邻域内的所有点，构建邻域点集的协方差矩阵，通过主成分分析求解协方差矩阵的特征值与特征向量，其中最小特征值对应的特征向量，即为该点的法向量。该方法计算简单，能够有效估计点云的法向量，但其估计结果受邻域大小的影响较大，同时对噪声与离群点较为敏感，需要在去噪操作完成后再进行法向量估计，以保证估计结果的精度。

\section{Transformer基础理论}

\subsection{自注意力机制}
自注意力机制是Transformer架构的核心组成部分，其核心作用是建模序列中各个元素之间的依赖关系，通过计算序列中每个元素与其他所有元素的关联权重，对特征进行加权聚合，从而捕捉序列中的长距离依赖关系。与CNN的卷积操作只能捕捉局部邻域内的特征关联不同，自注意力机制能够直接建模序列中任意两个元素之间的关联，无论两个元素在序列中的距离远近，都能有效捕捉其依赖关系，这也是Transformer相较于CNN与MLP最核心的优势。

缩放点积注意力是自注意力机制最常用的实现形式，其输入为三个矩阵：查询矩阵 $Q \in \mathbb{R}^{n \times d_k}$、键矩阵 $K \in \mathbb{R}^{m \times d_k}$、值矩阵 $V \in \mathbb{R}^{m \times d_v}$，其中 $n$ 为查询序列的长度，$m$ 为键-值序列的长度，$d_k$ 为查询与键向量的维度，$d_v$ 为值向量的维度。缩放点积注意力的计算过程可以表示为：
\begin{equation}
\text{Attention}(Q, K, V) = \text{softmax}\left( \frac{QK^T}{\sqrt{d_k}} \right) V
\end{equation}

在计算过程中，首先计算查询矩阵 $Q$ 与键矩阵 $K$ 的转置的点积，得到每个查询元素与所有键元素之间的相似度得分，得分的高低代表了两个元素之间的关联程度。随后除以 $\sqrt{d_k}$ 进行缩放，这一操作的核心目的是避免当 $d_k$ 较大时，点积的结果数值过大，导致softmax函数进入梯度饱和区，出现梯度消失的问题。再通过softmax函数将相似度得分转换为0到1之间的注意力权重，所有权重的和为1，权重的大小代表了对应值向量对当前查询向量的重要程度。最后将注意力权重与值矩阵 $V$ 进行加权求和，得到最终的注意力输出特征，该特征融合了序列中所有元素的信息，且根据元素之间的关联程度进行了加权。

为了进一步提升自注意力机制的特征建模能力，Transformer架构中采用了多头注意力机制，将查询、键、值矩阵通过多个独立的线性投影层，映射到多个不同的表示子空间中，在每个子空间中分别执行缩放点积注意力计算，再将多个头的输出结果进行拼接，通过线性投影层得到最终的输出。多头注意力机制的计算过程可以表示为：
\begin{equation}
\begin{split}
\text{MultiHead}(Q, K, V) &= \text{Concat}(\text{head}_1, \text{head}_2, ..., \text{head}_h) W^O \\
\text{where head}_i &= \text{Attention}(Q W_i^Q, K W_i^K, V W_i^V)
\end{split}
\end{equation}

其中 $h$ 为注意力头的数量，$W_i^Q \in \mathbb{R}^{d_{\text{model}} \times d_k}$、$W_i^K \in \mathbb{R}^{d_{\text{model}} \times d_k}$、$W_i^V \in \mathbb{R}^{d_{\text{model}} \times d_v}$ 为每个头对应的线性投影矩阵，$W^O \in \mathbb{R}^{h d_v \times d_{\text{model}}}$ 为最终的线性投影矩阵，$d_{\text{model}}$ 为Transformer模型的整体特征维度。在实际应用中，通常设置 $d_k = d_v = d_{\text{model}} / h$，保证每个注意力头的计算量与单头注意力相当。多头注意力机制能够让模型在不同的表示子空间中学习不同的特征关联，捕捉序列中更加丰富的依赖关系，相较于单头注意力具备更强的特征建模能力，能够更好地适配复杂的序列建模任务。

\subsection{编码器-解码器结构}
Transformer架构采用了经典的编码器-解码器结构，其中编码器负责对输入序列进行特征提取，得到包含全局上下文信息的序列特征；解码器负责基于编码器输出的特征，自回归地生成目标序列。在点云补全任务中，输入的残缺点云序列对应编码器的输入，需要生成的完整点云序列对应解码器的输出，因此编码器-解码器结构是点云补全网络的核心基础框架。

Transformer的编码器由多个相同的编码器层堆叠而成，每个编码器层主要包含两个子层：多头自注意力层与前馈网络层。在每个子层的前后，都采用了残差连接与层归一化操作，即每个子层的输出可以表示为 $\text{LayerNorm}(x + \text{Sublayer}(x))$，其中 $x$ 为子层的输入，$\text{Sublayer}(x)$ 为子层的计算结果。这种残差连接与层归一化的设计，能够有效缓解深层网络训练过程中的梯度消失问题，提升模型的训练稳定性与收敛速度，使得深层的Transformer架构能够被有效训练。

编码器层中的前馈网络层由两个线性变换层与一个非线性激活函数组成，其计算过程可以表示为：
\begin{equation}
\text{FFN}(x) = \max(0, x W_1 + b_1) W_2 + b_2
\end{equation}

其中 $W_1, W_2$ 为线性变换的权重矩阵，$b_1, b_2$ 为偏置项，通常在两个线性变换层之间采用ReLU（Rectified Linear Unit）或者GELU（Gaussian Error Linear Unit）作为激活函数。前馈网络层的核心作用是对注意力层输出的特征进行进一步的非线性变换，提升特征的表达能力，同时对每个位置的特征进行独立的处理，不改变序列的长度与维度。编码器的每一层都能够对输入序列的特征进行进一步的聚合与优化，通过多层编码器的堆叠，最终得到包含输入序列全局上下文信息的特征表示，为解码器的生成过程提供支撑。

Transformer的解码器同样由多个相同的解码器层堆叠而成，每个解码器层包含三个子层：掩码多头自注意力层、编码器-解码器注意力层与前馈网络层。与编码器层相同，每个子层也都采用了残差连接与层归一化操作，保证深层网络的训练稳定性。

掩码多头自注意力层的核心作用是保证序列生成的因果性，即在生成当前位置的元素时，只能利用已经生成的前面位置的元素，无法看到后面位置的元素，避免模型利用未来的信息进行作弊。因此需要对注意力矩阵进行掩码操作，将当前位置之后的元素的注意力权重设置为负无穷，经过softmax函数后权重变为0，从而实现对未来信息的屏蔽，保证序列生成的因果合理性。

编码器-解码器注意力层也被称为交叉注意力层，其查询矩阵来自于前一个解码器层的输出，键矩阵与值矩阵来自于编码器的最终输出。该层的核心作用是让解码器在生成目标序列的每一个位置时，都能够关注到输入序列中对应的相关特征，建立输入序列与目标序列之间的关联。这对于点云补全任务至关重要，能够让解码器在生成完整点云的过程中，充分利用输入残缺点云的几何特征，保证生成的完整点云与输入的残缺点云在结构上的一致性，避免出现结构失真的问题。解码器的前馈网络层与编码器中的结构完全相同，用于对特征进行进一步的非线性变换，提升特征的表达能力。

\subsection{位置编码}
自注意力机制的核心特性之一是排列不变性，即对于输入序列，任意改变序列中元素的排列顺序，自注意力机制的输出结果不会发生改变。这一特性虽然适配了点云的无序性，但也导致自注意力机制无法捕捉序列中元素的位置信息，而对于序列处理任务而言，元素的位置信息往往包含了重要的语义信息，对于点云处理任务而言，点的三维空间位置更是核心的几何信息。为了解决这一问题，Transformer架构中引入了位置编码，将序列中元素的位置信息注入到输入特征中，让模型能够捕捉到元素的位置关联。

位置编码的核心目标是为序列中的每个位置生成一个唯一的位置编码向量，该向量与输入特征的维度相同，可以直接与输入特征进行相加，作为编码器的输入。常用的位置编码主要分为绝对位置编码与相对位置编码两大类，其中正弦余弦位置编码是最经典的绝对位置编码方法，其通过正弦与余弦函数为每个位置生成固定的位置编码，无需学习，计算过程可以表示为：
\begin{equation}
\begin{split}
PE_{(pos, 2i)} &= \sin\left( pos / 10000^{2i / d_{\text{model}}} \right) \\
PE_{(pos, 2i+1)} &= \cos\left( pos / 10000^{2i / d_{\text{model}}} \right)
\end{split}
\end{equation}

其中 $pos$ 为元素在序列中的位置索引，$i$ 为位置编码向量的维度索引，$d_{\text{model}}$ 为特征的维度。正弦余弦位置编码的优势在于，其能够通过三角函数的性质，捕捉序列中不同位置之间的相对位置关系，对于任意两个位置 $pos+k$ 与 $pos$，其位置编码之间的关系可以通过线性变换表示，能够让模型轻松学习到序列的相对位置信息。同时，对于训练过程中没有出现过的序列长度，也能够生成有效的位置编码，具备良好的泛化能力。

但对于点云处理任务而言，传统的一维位置编码无法直接适用，因为点云是三维空间中的离散点，不存在固定的一维序列位置，其位置信息体现在三维空间坐标中，而非序列中的索引位置。因此在点云Transformer架构中，通常采用基于点三维坐标的几何位置编码，直接将点的三维坐标通过线性变换映射到与特征相同的维度，作为位置编码与点特征相加，或者将点的三维坐标直接融入到注意力权重的计算中，生成几何感知的注意力权重，让模型能够直接捕捉到点云中点与点之间的三维空间位置关联，这也是点云Transformer与二维图像、自然语言处理中Transformer的核心区别之一\cite{yu2021pointr}。

\section{点云补全核心技术}

\subsection{编码器-解码器补全框架}
编码器-解码器框架是当前点云补全任务最主流的基础架构，几乎所有的深度学习点云补全方法都基于这一框架进行设计与优化，其核心思想是通过编码器将输入的残缺点云映射到高维特征空间，得到包含物体几何信息的特征表示，再通过解码器从高维特征中映射回三维点云空间，生成完整的点云。这一框架实现了端到端的点云补全，无需人工设计几何先验，能够通过数据驱动的方式，学习从残缺点云到完整点云的映射关系。

编码器是补全框架的核心，其核心目标是从输入的残缺点云中，提取能够有效表征物体整体几何形状与局部结构细节的特征，为后续的补全生成提供足够的信息。编码器的特征提取能力，直接决定了补全结果的上限，若编码器无法有效提取输入点云的核心几何特征，解码器也无法生成高质量的完整点云。早期的补全方法大多基于PointNet++架构设计编码器，通过多层的点特征提取与最大池化操作，得到输入点云的全局特征，但这种方法往往会丢失大量的局部细节信息，无法有效建模点之间的长距离依赖关系，对于大面积缺失的点云，难以提取有效的全局形状特征。基于Transformer的补全方法则采用Transformer编码器作为特征提取的核心，通过自注意力机制建模点云中点与点之间的全局关联，同时结合多尺度特征提取，同时捕捉物体的全局形状上下文与局部几何细节，得到更具表征能力的特征表示，为高质量的补全生成奠定基础\cite{yu2023adapointr}。

解码器是补全框架的生成核心，其核心目标是从编码器提取的特征中，生成完整的高密度点云，解码器的设计直接决定了补全结果的质量与细节保真度。早期的补全方法大多采用折叠操作作为解码器的核心，其核心思想是将二维平面上的规则网格点，通过多层感知机映射到三维空间，折叠成物体的三维形状。这种方法计算简单，能够快速生成完整的点云，但只能生成拓扑结构较为简单的物体表面，对于复杂拓扑结构的物体，无法有效还原其几何形状，同时容易出现局部细节丢失的问题。后续的研究提出了多种不同的解码器设计，包括点反卷积、分形点生成、多阶段上采样等，而基于Transformer的解码器则通过自注意力机制与交叉注意力机制，建模生成点之间的关联，同时充分利用编码器提取的多尺度特征，实现更精准的点云生成，能够更好地还原物体的复杂拓扑结构与精细几何细节\cite{xiang2021snowflakenet}。

\subsection{常用损失函数}
损失函数是点云补全网络训练的核心，其定义了网络的优化目标，直接决定了网络的收敛方向与最终的补全效果。点云补全任务的核心目标是让生成的完整点云与真实的完整点云在几何结构、点分布、细节特征上尽可能一致，因此需要设计对应的损失函数，从多个维度对生成点云进行约束。常用的损失函数主要包括以下几类。

倒角距离（Chamfer Distance, CD）是点云补全任务中最基础、最常用的损失函数，其核心思想是计算两个点云之间的双向平均最小距离，衡量生成点云与真实点云之间的整体几何差异。对于生成点云 $\mathcal{P}$ 与真实点云 $\mathcal{Q}$，倒角距离的计算过程可以表示为：
\begin{equation}
\text{CD}(\mathcal{P}, \mathcal{Q}) = \frac{1}{|\mathcal{P}|} \sum_{p \in \mathcal{P}} \min_{q \in \mathcal{Q}} \| p - q \|_2^2 + \frac{1}{|\mathcal{Q}|} \sum_{q \in \mathcal{Q}} \min_{p \in \mathcal{P}} \| q - p \|_2^2
\end{equation}

其中 $|\mathcal{P}|$ 与 $|\mathcal{Q}|$ 分别为生成点云与真实点云的点数。倒角距离的前半部分计算了生成点云中的每个点到真实点云的最近点距离的平均值，约束生成点云的所有点都尽可能贴近真实点云的表面，避免生成离群点；后半部分计算了真实点云中的每个点到生成点云的最近点距离的平均值，约束生成点云能够完整覆盖真实点云的所有表面区域，避免出现补全不全的问题。倒角距离的计算效率高，对两个点云的点数没有严格的匹配要求，能够有效衡量两个点云之间的整体几何差异，是点云补全任务的核心损失函数。但倒角距离也存在明显的不足，其只关注点与点之间的最小距离，无法约束生成点云的分布均匀性，容易导致生成的点云出现局部聚集、密度不均的问题，影响后续的三维重建效果。

推土机距离（Earth Mover's Distance, EMD）是另一种常用的点云相似性度量损失函数，其核心思想是将两个点云看作两个概率分布，计算将一个分布转换为另一个分布所需的最小累积成本，能够更好地衡量两个点云之间的分布匹配程度。对于生成点云 $\mathcal{P}$ 与真实点云 $\mathcal{Q}$，当两个点云的点数相同时，推土机距离的计算过程可以表示为：
\begin{equation}
\text{EMD}(\mathcal{P}, \mathcal{Q}) = \min_{\phi: \mathcal{P} \to \mathcal{Q}} \frac{1}{|\mathcal{P}|} \sum_{p \in \mathcal{P}} \| p - \phi(p) \|_2
\end{equation}

其中 $\phi$ 为从生成点云 $\mathcal{P}$ 到真实点云 $\mathcal{Q}$ 的双射映射，即每个生成点都对应唯一的真实点，保证了两个点云之间的最优匹配。推土机距离能够找到两个点云之间的全局最优匹配，更好地衡量两个点云之间的整体分布差异，相较于倒角距离，其对生成点云的分布均匀性有更好的约束效果，生成的点云更加均匀，细节质量更高。但推土机距离的计算复杂度远高于倒角距离，尤其是在点云点数较多时，计算成本极高，无法应用于高分辨率点云的训练。因此在实际训练中，通常只在低分辨率的粗生成阶段采用推土机距离，而在高分辨率的细生成阶段采用倒角距离，平衡计算成本与约束效果。

除了上述两种基础的几何损失函数之外，为了进一步提升生成点云的质量，学者们提出了多种辅助损失函数，从不同维度对生成点云进行约束。密度感知损失的核心目标是约束生成点云的局部密度，让生成点云的密度分布与真实点云保持一致，解决倒角距离带来的密度不均问题。其核心原理是对每个点，计算其指定邻域内的点数量，得到局部密度信息，通过损失函数约束生成点云与真实点云的局部密度差异，保证生成点云的分布均匀性，为后续的三维重建提供高质量的点云数据。排斥损失的核心目标是避免生成的点出现聚集重叠的问题，其核心原理是对每个生成点，计算其与邻域内其他生成点之间的距离，通过损失函数惩罚距离过近的点，让生成的点之间保持合理的距离，避免出现点云聚集的问题，提升点云的整体质量。

\section{三维重建基础方法}
点云补全的最终目标是为后续的三维重建任务提供高质量的完整点云，从而得到物体的三维网格模型，实现物体的数字化表示。三维表面重建的核心任务是从离散的三维点云中，重建出物体的连续表面网格模型，将离散的点云转换为连续的、可编辑的三维数字模型。本文中的主要三位表面重建算法为泊松重建。% 常用的重建方法主要包括泊松重建与基于球旋转算法的重建，这两类方法也是实际应用中最主流的三维表面重建方法。

泊松重建（Poisson Surface Reconstruction）是一种基于隐函数的三维表面重建方法，其核心思想是通过求解泊松方程，得到三维空间中的指示函数，将物体的表面定义为指示函数的等值面，从而重建出物体的封闭、水密的网格模型。泊松重建的核心优势在于，其能够利用点云的法向量信息，构建全局的隐函数，对噪声与离群点具有较强的鲁棒性，能够重建出光滑、连续的表面网格，是当前高质量三维重建最常用的方法之一。

泊松重建的核心原理基于向量场的散度定理，其首先将输入点云的法向量场，转换为三维空间中的梯度向量场，而指示函数的梯度应该与该梯度向量场尽可能一致。指示函数是一个定义在三维空间中的标量函数，在物体内部的取值为1，在物体外部的取值为0，其梯度在物体的表面位置会出现突变，因此指示函数的等值面就对应了物体的表面。根据梯度向量场求解指示函数的过程，最终可以转化为求解一个泊松方程，通过有限元方法求解该泊松方程，即可得到三维空间中的指示函数。在得到指示函数后，通过移动立方体（Marching Cubes）算法，提取指示函数的0.5等值面，即可得到物体的三维三角形网格模型。泊松重建的细节原理将在第\ref{cha:reconstruct}章中展开论述，在此不做过多赘述。

泊松重建的效果受输入点云的法向量精度、重建深度参数的影响较大，重建深度决定了重建网格的分辨率，深度越大，重建的网格越精细，细节越丰富，但同时计算量也会越大，对噪声也会更加敏感。泊松重建能够生成水密的网格模型，不存在孔洞与缝隙，非常适合后续的3D打印、有限元分析、数字孪生等应用。但对于存在大面积缺失的点云，泊松重建无法准确还原物体的表面结构，会出现严重的拓扑结构失真，这也是需要先进行点云补全的核心原因。只有输入完整、均匀的高质量点云，泊松重建才能得到高精度的三维网格模型。

% \subsection{基于球旋转算法的三维表面重建}
% 球旋转算法（Ball Pivoting Algorithm, BPA）是一种基于空间几何的直接三维表面重建方法，其核心思想是通过一个固定半径的球，在输入的稠密点云上滚动，当球同时碰到三个不共线的点时，这三个点就构成一个三角形面片，同时这三个点形成的三角形边作为球滚动的旋转轴，球沿着该轴继续滚动，碰到下一个点时，生成新的三角形面片。重复这一过程，直到球无法再滚动到新的点，最终生成覆盖整个点云的三角形网格模型。

% 球旋转算法的计算过程简单，计算效率远高于泊松重建，能够快速从稠密点云中生成表面网格，同时生成的网格顶点直接采用原始点云中的点，不会改变原始点云的几何信息，能够精准还原点云的几何结构，保留原始点云的所有细节特征。但球旋转算法也存在明显的局限性，其对输入点云的密度与均匀性要求极高，对于稀疏、非均匀的点云，固定半径的球无法同时碰到三个点，会导致重建的表面出现大量的孔洞；同时对于点云中的噪声与离群点，球旋转算法非常敏感，会生成大量错误的三角形面片，影响重建效果。因此球旋转算法通常只适用于分布均匀、密度较高、噪声较少的理想点云，对于实际采集的残缺、稀疏点云，需要先进行补全与预处理，提升点云的质量，才能得到较好的重建效果。

\section{本章小结}
本章系统阐述了本文研究涉及的相关理论与基础技术，为后续的算法设计、实验验证与系统实现奠定了坚实的理论基础。首先介绍了点云数据的基础定义与核心特性，以及点云预处理过程中常用的去噪、下采样、归一化、法向量估计等操作，明确了点云处理任务的基础要求与核心难点。随后详细阐述了Transformer架构的核心理论，包括自注意力机制、多头注意力机制、编码器-解码器结构与位置编码方法，分析了Transformer在长距离依赖建模上的核心优势，以及在点云处理任务中的适配方法。在此基础上，介绍了点云补全任务的核心技术，包括编码器-解码器补全框架、由粗到精的生成策略，以及常用的损失函数，分析了不同技术的核心原理与优缺点。最后简单介绍了本文用于三维表面重建的泊松重建法，明确了点云补全与三维重建之间的关联，以及点云补全对三维重建效果的重要影响。本章介绍的相关理论与技术，是本文后续算法设计的核心依据，也为实验结果的分析提供了理论支撑。
